\part{Physical System}
\chapter{Physical Constraints}
\label{sec:physical-constraints}
Physical constraints are specified using the \texttt{phys.} prefix and control the physical layout and manufacturing requirements of the PCB.

\section{Units and SI Prefixes}
\textbf{Warning:} All physical values use base SI units (meters, ohms, farads, etc.). A value like \texttt{min=15} means 15 \textbf{meters}, not 15 millimeters!

Always use SI prefixes for typical PCB dimensions:
\begin{itemize}
\item \texttt{min=15m} -- 15 millimeters ($15 \times 10^{-3}$ m)
\item \texttt{min=15u} -- 15 micrometers ($15 \times 10^{-6}$ m)
\item \texttt{d=0.25m} -- 0.25 millimeters (250 $\mu$m)
\end{itemize}

SI prefixes commonly used in PCB design: \texttt{n} ($10^{-9}$), \texttt{u} ($10^{-6}$), \texttt{m} ($10^{-3}$), \texttt{k} ($10^{3}$), \texttt{M} ($10^{6}$), \texttt{G} ($10^{9}$). All SI prefixes (quecto through quetta) are supported — see Section~\ref{sec:generic-types} for the full table.

\section{Priority System}
All physical constraints accept a priority parameter that determines how conflicts are resolved:
\begin{verbatim}
Priorities:
  (none) or -d     Default
  -nth             Nice to have
  -imp             Important
  -vimp            Very important
  -oeo             Only explicit overrides
  -eo              Explicit override
  -ne = -nonneg    No exceptions

Priority Chain: 
  -d < -nth < -imp < -vimp < -oeo < -eo < -ne = -nonneg
\end{verbatim}

\section{Proximity Constraint}
Controls the distance between sets of objects:
\begin{verbatim}
phys.prox <set_a> <set_b> 
    min=<min_distance> max=<max_distance> 
    <priority>
phys.prox <set> minim <priority>
\end{verbatim}
\begin{itemize}
\item \texttt{min} -- Minimum distance between edges of set a and set b
\item \texttt{max} -- Maximum distance anything in set a can be from anything in set b
\item \texttt{minim} -- Tries to minimize distance between all objects in the set
\end{itemize}

Example:
\begin{verbatim}
phys.prox $hvsys $lvsys min=15m -ne
phys.prox $noisy_loop minim -vimp
\end{verbatim}

\section{Creepage Constraint}
Controls the surface distance between sets (affected by slots, vias, traces, and conductive material):
\begin{verbatim}
phys.creep <set_a> <set_b> 
    d=<min_creepage_distance> <priority>
\end{verbatim}

Note: Unlike \texttt{phys.prox} (Euclidean distance), \texttt{phys.creep} measures surface distance. Conductive material contributes 0 distance.

Example:
\begin{verbatim}
phys.creep $hvsys $lvsys d=30m -nonneg
\end{verbatim}

\section{Layer Constraint}
Assigns objects to specific board layers:
\begin{verbatim}
phys.layer <set> <layer_number> <priority>
\end{verbatim}
Layer 0 is the top layer; the bottom layer is layer $n$ (where $n$ is the total number of layers).

Example:
\begin{verbatim}
phys.layer *p 0 -oeo
\end{verbatim}

\section{Board Definition}
Defines the PCB stackup characteristics. Can only be invoked once; a \texttt{Runtime Error} is raised if called a second time.
\begin{verbatim}
phys.board N=<layers> T=[<thicknesses>] 
           W=[<weights>] M=[<materials>]
\end{verbatim}
\begin{itemize}
\item \texttt{N} -- Number of board layers
\item \texttt{T} -- List of thicknesses between layers in meters (length $n-1$)
\item \texttt{W} -- List of copper weights in oz/ft$^2$ (length $n$)
\item \texttt{M} -- List: [dielectric, conductor, surface-plating]
\end{itemize}

Default: \texttt{phys.board N=2 T=[0.0016] W=[1, 1] 
         M=["FR4", "copper", "HASL"]}

Example:
\begin{verbatim}
phys.board N=2 T=[0.0016] W=[2, 2] 
           M=["FR4", "copper", "ENIG"]
\end{verbatim}

\section{Footprint Unification}
Merges multiple parts into a single physical unit on the PCB. The parts share one combined footprint and reference designator, but retain independent electrical pins. In the schematic, united parts are grouped together.
\begin{verbatim}
phys.unite { <part>, <part>, ... }
\end{verbatim}

\begin{itemize}
\item The argument is a set of part references.
\item United parts share a single footprint on the PCB.
\item Reference designators are combined (e.g., \texttt{R1.R2.R3}).
\item Pins remain independently accessible for net connections.
\end{itemize}

Example:
\begin{verbatim}
phys.unite { R1, R2, R3 }
#c# R1, R2, R3 share one footprint on the PCB
#c# Schematic groups them together
#c# R1[1], R2[1], R3[1] are still individually accessible
\end{verbatim}

\section{Net Shrink}
Minimizes the physical extent of a net:
\begin{verbatim}
phys.shrinknet <net> <priority>
\end{verbatim}

A \texttt{Value Error} is raised if the referenced net does not exist.

Example:
\begin{verbatim}
elec.net "Vsw"
phys.shrinknet Vsw -vimp
\end{verbatim}

\section{Plane Creation}
Creates a copper plane on a layer connected to a net:
\begin{verbatim}
phys.mkplane <layer_number> <net>
\end{verbatim}

A \texttt{Value Error} is raised if the referenced net or layer does not exist.

Example:
\begin{verbatim}
elec.net "Gnd"
phys.mkplane 1 Gnd
\end{verbatim}

\section{Text Annotation}
Places text on the board as a physical object:
\begin{verbatim}
phys.mktext txt = <text> 
    font = <font_name> 
    size = <point_size> 
    material = <material> 
    layer = <layer>
\end{verbatim}
\begin{itemize}
\item \texttt{txt} -- Text string to display
\item \texttt{font} -- Font name (e.g., \texttt{"Ubuntu Mono"})
\item \texttt{size} -- Font size in points
\item \texttt{material} -- Material (\texttt{"cu"} for copper silkscreen)
\item \texttt{layer} -- Target layer (e.g., \texttt{T} for top silkscreen). Note: In \texttt{phys.mktext}, \texttt{T} refers to the top silkscreen layer, not the top copper layer as in \texttt{phys.layer}.
\end{itemize}

Example:
\begin{verbatim}
phys.mktext txt = "HV" font = "Ubuntu Mono" 
    size = 9 material = "cu" layer = T
\end{verbatim}

\section{Via Fill}
Fills a pad with vias for thermal management:
\begin{verbatim}
phys.viafill <pad> v=<via_hole_diameter> 
    d=<via_spacing> [p=<pattern>] [l=<layers>]
\end{verbatim}
\begin{itemize}
\item \texttt{v} -- Via hole diameter (use \texttt{m} for mm)
\item \texttt{d} -- Minimum spacing between via centers
\item \texttt{p} -- Fill pattern (default: \texttt{grid})
\item \texttt{l} -- Layer span (default: \texttt{all}). Can be \texttt{all} or a pair like \texttt{1-2}
\end{itemize}

\paragraph{Edge Cases}
\begin{itemize}
\item A \texttt{Value Error} is raised if \texttt{v} or \texttt{d} is $\leq 0$.
\item A \texttt{Value Error} is raised if \texttt{d} $<$ \texttt{v} (vias would overlap).
\item A \texttt{Value Error} is raised if an invalid pattern is specified.
\item A \texttt{Value Error} is raised if \texttt{l} references layers outside the board stackup.
\item A \texttt{Type Error} is raised if the pad argument is not a valid pin reference.
\end{itemize}

\subsection{Pattern Types}

All patterns tile the pad area with vias at spacing \texttt{d}. Each via is placed at the vertices of the respective tiling.

\paragraph{grid}
Rectangular lattice. Vias at $(i \cdot d, \; j \cdot d)$ for integers $i, j$.
\begin{verbatim}
phys.viafill Q1[d] v=0.25m d=0.5m p=grid
\end{verbatim}

\paragraph{hex}
Densest circle packing arrangement. Each via has 6 nearest neighbors at distance $d$. Approximately 15\% denser than grid, best thermal performance due to maximized copper coverage.
\begin{verbatim}
phys.viafill Q1[d] v=0.25m d=0.5m p=hex
\end{verbatim}

\paragraph{stagger}
Offset rows forming the same spacing as hex but with alternating row offsets. Equivalent to hex for thermal purposes but creates straighter routing channels between rows.
\begin{verbatim}
phys.viafill Q1[d] v=0.25m d=0.5m p=stagger
\end{verbatim}

\paragraph{radial}
Concentric rings emanating from pad center. Ring radii at $r_n = d \cdot n$ for integer $n \geq 1$. Optimal for round thermal pads under ICs.
\begin{verbatim}
phys.viafill U1[pin5] v=0.3m d=0.8m p=radial
\end{verbatim}

\paragraph{random}
Poisson-disk distributed vias with minimum spacing $d$. Avoids periodic structures that can cause resonance or EMI.
\begin{verbatim}
phys.viafill Q1[d] v=0.25m d=0.5m p=random
\end{verbatim}

\subsection{Examples}
\begin{verbatim}
#c# Standard grid fill for MOSFET drain
phys.viafill Q1[d] v=0.25m d=0.5m p=grid

#c# Dense hex fill for high-current pad
phys.viafill Q1[d] v=0.2m d=0.4m p=hex

#c# Radial fill for IC thermal pad
phys.viafill U1[pin5] v=0.3m d=0.8m p=radial

#c# Random fill to avoid EMI resonance
phys.viafill Q1[d] v=0.25m d=0.5m p=random
\end{verbatim}

\section{Connection Current}
Sets the minimum RMS current a connection must carry:
\begin{verbatim}
phys.connectioncurrent <connection | set> 
    <rms_current> <priority>
\end{verbatim}

Example:
\begin{verbatim}
phys.connectioncurrent ( J1[2] | F1[1] ) 30 -oeo
\end{verbatim}

\section{Alignment}
Aligns parts along an axis with specified spacing:
\begin{verbatim}
phys.align <list_of_parts> 
    axis=<x|y> d=<edge_spacing> <priority>
\end{verbatim}

Example:
\begin{verbatim}
phys.align [ Q1, Q2 ] axis=x d=5m -vimp
\end{verbatim}

\section{Connection Matching}
Minimizes variance of a property across a set of connections:
\begin{verbatim}
phys.connmatch <property> 
    <set_of_connections> <priority>
\end{verbatim}
The set of connections is a \texttt{set} of connection objects (e.g., connections between parts and nets), or a variable holding such a set.

\subsection{Properties}
\begin{itemize}
\item \texttt{l} -- Physical trace length
\item \texttt{res} -- DC resistance
\item \texttt{L} -- Self-inductance
\item \texttt{cap} -- Capacitance to reference plane
\item \texttt{imp} -- Characteristic impedance
\item \texttt{t} -- Signal propagation delay
\end{itemize}

\subsection{Value Calculation Methods}

\subsubsection{Calculation Mode Selection}
The calculation method can be selected using the \texttt{syst.cm.calc} flag:
\begin{verbatim}
syst.cm.calc <mode>
\end{verbatim}

Available modes:
\begin{itemize}
\item \texttt{approx} -- (Default) Uses simplified closed-form formulas for fast computation. Suitable for most designs.
\item \texttt{calculus} -- Uses numerical integration and differential equations for higher accuracy. Computationally more expensive.
\item \texttt{hybrid} -- Uses closed-form formulas for initial estimates, then refines with calculus-based methods for critical nets.
\end{itemize}

Example:
\begin{verbatim}
syst.cm.calc calculus
phys.connmatch imp $high_speed_lines -vimp
\end{verbatim}

\subsubsection{Length (len)}
The physical length of the trace path from source pin to destination pin, measured in millimeters. Calculated as the sum of all trace segments.

\subsubsection{Resistance (res)}
DC resistance is calculated using:
\begin{equation}
R = \frac{\rho \cdot l}{W \cdot T}
\end{equation}
where:
\begin{itemize}
\item $\rho$ -- Resistivity of the conductor (copper: $1.72 \times 10^{-8}$ $\Omega\cdot$m)
\item $l$ -- Trace length (mm)
\item $W$ -- Trace width (mm)
\item $T$ -- Trace thickness (mm, derived from copper weight in oz/ft$^2$)
\end{itemize}

\paragraph{Calculus Mode}
When \texttt{syst.cm.calc calculus} is active, skin effect is considered:
\begin{equation}
R_{ac} = R_{dc} \cdot \left(1 + \frac{1}{12} \cdot \left(\frac{\delta}{T}\right)^4\right)
\end{equation}
where $\delta = \sqrt{\frac{2\rho}{\omega\mu}}$ is the skin depth at frequency $\omega$.

\subsubsection{Inductance (ind)}
Self-inductance is approximated using the loop inductance formula for a microstrip:
\begin{equation}
L \approx \frac{\mu_0 \cdot l_{trace}}{2\pi} \cdot \ln\left(\frac{2h}{w + t}\right)
\end{equation}
where:
\begin{itemize}
\item $\mu_0$ -- Permeability of free space ($4\pi \times 10^{-7}$ H/m)
\item $l_{trace}$ -- Trace length
\item $h$ -- Height above reference plane
\item $w$ -- Trace width
\item $t$ -- Trace thickness
\end{itemize}

\paragraph{Calculus Mode}
When \texttt{syst.cm.calc calculus} is active, mutual inductance and frequency-dependent effects are computed using numerical integration:
\begin{equation}
L_{total} = \int_0^{l_{trace}} \frac{\mu_0}{2\pi} \cdot \ln\left(\frac{2h(x)}{w(x) + t}\right) dx
\end{equation}
This accounts for trace width variations and non-uniform height above the reference plane.

\subsubsection{Capacitance (cap)}
Capacitance to the reference plane is calculated using the microstrip capacitance formula:
\begin{equation}
C \approx \frac{\varepsilon_0 \cdot \varepsilon_r \cdot w \cdot l_{trace}}{h}
\end{equation}
where:
\begin{itemize}
\item $\varepsilon_0$ -- Permittivity of free space ($8.854 \times 10^{-12}$ F/m)
\item $\varepsilon_r$ -- Relative permittivity of the dielectric (FR4: $\approx 4.5$)
\item $w$ -- Trace width
\item $l_{trace}$ -- Trace length
\item $h$ -- Height above reference plane
\end{itemize}

\paragraph{Calculus Mode}
When \texttt{syst.cm.calc calculus} is active, fringing effects and dielectric variations are included:
\begin{equation}
C = \int_0^{l_{trace}} \frac{\varepsilon_0 \cdot \varepsilon_r(x) \cdot w_{eff}(x)}{h(x)} dx
\end{equation}
where $w_{eff}$ includes fringing capacitance contributions and $\varepsilon_r(x)$ accounts for material variations along the trace.

\subsubsection{Impedance (imp)}
Characteristic impedance for a microstrip is calculated as:
\begin{equation}
Z_0 \approx \frac{87}{\sqrt{\varepsilon_r + 1.41}} \cdot \ln\left(\frac{5.98h}{0.8w + t}\right)
\end{equation}
For stripline:
\begin{equation}
Z_0 \approx \frac{60}{\sqrt{\varepsilon_r}} \cdot \ln\left(\frac{4b}{0.67(0.8w + t)}\right)
\end{equation}
where $b$ is the total dielectric thickness for stripline configurations.

\paragraph{Calculus Mode}
When \texttt{syst.cm.calc calculus} is active, the full telegrapher's equations are solved numerically:
\begin{equation}
Z_0 = \sqrt{\frac{R + j\omega L}{G + j\omega C}}
\end{equation}
where $R$, $L$, $G$, and $C$ are per-unit-length parameters computed from the trace geometry and material properties. This accounts for frequency-dependent losses and dispersion.

\subsubsection{Propagation Time (t)}
Signal propagation delay is calculated as:
\begin{equation}
t = \frac{l_{trace}}{v_p} = \frac{l_{trace} \cdot \sqrt{\varepsilon_{eff}}}{c}
\end{equation}
where:
\begin{itemize}
\item $v_p$ -- Propagation velocity
\item $c$ -- Speed of light in vacuum ($3 \times 10^8$ m/s)
\item $\varepsilon_{eff}$ -- Effective permittivity (depends on geometry and dielectric)
\end{itemize}

\paragraph{Calculus Mode}
When \texttt{syst.cm.calc calculus} is active, propagation delay is computed by integrating along the trace path:
\begin{equation}
t = \int_0^{l_{trace}} \frac{\sqrt{\varepsilon_{eff}(x)}}{c} dx
\end{equation}
This accounts for varying dielectric thickness, material transitions, and impedance discontinuities along the trace.

\subsection{Matching Objective}
The solver minimizes the scaled variance of the selected property across all connections in the set:
\begin{equation}
\text{minimize} \quad \frac{\text{var}(P)}{||S||}
\end{equation}
where $P$ is the property value for each connection and $||S||$ is the cardinality of the connection set.

\subsection{Example}
\begin{verbatim}
phys.connmatch t $CPU_to_RAM_lines -vimp
\end{verbatim}
This ensures all CPU-to-RAM signal lines have matched propagation delays, critical for parallel data buses.


\chapter{Solver Model}
\label{sec:solver-model}

\section{Overview}
The EDADSL solver transforms declarative physical constraints into an optimized PCB layout. The process involves collecting constraints, selecting an optimization algorithm, and generating the final layout through placement and routing phases.

\section{Algorithm Selection}
The solver algorithm can be selected using the \texttt{syst.alg} flag:
\begin{verbatim}
syst.alg <algorithm_name>
\end{verbatim}

\subsection{Available Algorithms}

\subsubsection{Simulated Annealing (sa)}
A probabilistic optimization technique inspired by the annealing process in metallurgy. The algorithm explores the solution space by accepting worse solutions with a probability that decreases over time.

\paragraph{Objective Function}
Let $\mathbf{x} \in \mathcal{X}$ represent the complete layout state (component positions, trace routes, via placements). The objective function $f: \mathcal{X} \rightarrow \mathbb{R}$ quantifies the total constraint violation cost:
\begin{equation}
f(\mathbf{x}) = \sum_{i=1}^{n} w_i \cdot v_i(\mathbf{x})
\end{equation}
where $w_i$ is the priority weight of constraint $i$ and $v_i(\mathbf{x})$ is the violation measure (0 if satisfied, positive otherwise).

\paragraph{Neighbourhood Structure}
A neighbour state $\mathbf{x}'$ is generated by applying a random perturbation $\delta$ to the current state:
\begin{equation}
\mathbf{x}' = \mathbf{x} + \delta, \quad \delta \sim \mathcal{U}(-\epsilon, \epsilon)
\end{equation}
where $\epsilon$ is the maximum perturbation magnitude, adaptively scaled based on acceptance history.

\paragraph{Acceptance Probability}
The Metropolis acceptance criterion determines whether to move to a neighbour state:
\begin{equation}
P(\text{accept}) = \begin{cases}
1 & \text{if } f(\mathbf{x}') \leq f(\mathbf{x}) \\
\exp\left(-\frac{f(\mathbf{x}') - f(\mathbf{x})}{T}\right) & \text{if } f(\mathbf{x}') > f(\mathbf{x})
\end{cases}
\end{equation}
where $T > 0$ is the current temperature parameter.

\paragraph{Cooling Schedule}
The temperature decreases according to a geometric cooling schedule:
\begin{equation}
T_{k+1} = \alpha \cdot T_k, \quad \alpha \in (0, 1)
\end{equation}
with typical values $\alpha \in [0.95, 0.999]$. The initial temperature $T_0$ is set such that the initial acceptance probability for a worst-case perturbation is approximately $0.8$:
\begin{equation}
T_0 = -\frac{\Delta f_{\max}}{\ln(0.8)}
\end{equation}

\paragraph{Convergence Criterion}
The algorithm terminates when one of the following conditions is met:
\begin{itemize}
\item Temperature threshold: $T_k < T_{\min}$ (typically $10^{-6}$)
\item No improvement: $f(\mathbf{x}_{\text{best}})$ has not improved for $N_{\text{stale}}$ iterations
\item Maximum iterations: $k > k_{\max}$
\end{itemize}

\paragraph{Algorithm}
\begin{enumerate}
\item Initialize $\mathbf{x}_0$ randomly, set $T_0$, $k = 0$
\item Generate neighbour $\mathbf{x}'$ from $\mathbf{x}_k$
\item Compute $\Delta f = f(\mathbf{x}') - f(\mathbf{x}_k)$
\item If $\Delta f \leq 0$ or $\mathcal{U}(0,1) < \exp(-\Delta f / T_k)$, set $\mathbf{x}_{k+1} = \mathbf{x}'$
\item Otherwise, set $\mathbf{x}_{k+1} = \mathbf{x}_k$
\item Update $T_{k+1} = \alpha \cdot T_k$
\item If convergence not met, $k \leftarrow k+1$, goto 2
\end{enumerate}

\paragraph{Properties}
\begin{itemize}
\item Time complexity: $O(k_{\max} \cdot |\mathcal{N}(\mathbf{x})|)$ where $|\mathcal{N}(\mathbf{x})|$ is the neighbourhood evaluation cost
\item Probability of finding global optimum approaches 1 as $k_{\max} \rightarrow \infty$ with sufficiently slow cooling
\item Space complexity: $O(|\mathbf{x}|)$ for storing the current and best states
\end{itemize}

\subsubsection{Genetic Algorithm (ga)}
An evolutionary optimization method that maintains a population of candidate solutions and evolves them through biologically-inspired operators: selection, crossover, and mutation.

\paragraph{Population Representation}
The population at generation $t$ is a set of $N$ chromosomes:
\begin{equation}
P(t) = \{\mathbf{x}_1^{(t)}, \mathbf{x}_2^{(t)}, \ldots, \mathbf{x}_N^{(t)}\}
\end{equation}
Each chromosome $\mathbf{x}_i^{(t)} \in \mathcal{X}$ encodes a complete layout state as a gene vector:
\begin{equation}
\mathbf{x}_i^{(t)} = (g_1, g_2, \ldots, g_m)
\end{equation}
where each gene $g_j$ represents a design variable (position, route choice, etc.).

\paragraph{Fitness Function}
The fitness of chromosome $\mathbf{x}_i$ is derived from the objective function:
\begin{equation}
\text{fit}(\mathbf{x}_i) = \frac{1}{1 + f(\mathbf{x}_i)}
\end{equation}
where $f(\mathbf{x}_i)$ is the constraint violation cost. Higher fitness indicates better solutions.

\paragraph{Selection}
Parent selection uses tournament selection of size $\tau$:
\begin{enumerate}
\item Sample $\tau$ individuals uniformly from $P(t)$
\item Select the individual with highest fitness as parent
\end{enumerate}
This is repeated to select $N$ parents for mating.

\paragraph{Crossover}
Single-point crossover with probability $p_c$:
\begin{equation}
\text{crossover}(\mathbf{x}_i, \mathbf{x}_j) = \begin{cases}
(g_1^{(i)}, \ldots, g_c^{(i)}, g_{c+1}^{(j)}, \ldots, g_m^{(j)}) & \text{if } \mathcal{U}(0,1) < p_c \\
\mathbf{x}_i & \text{otherwise}
\end{cases}
\end{equation}
where $c \sim \mathcal{U}\{1, m-1\}$ is the crossover point.

\paragraph{Mutation}
Per-gene mutation with probability $p_m$:
\begin{equation}
g_j' = \begin{cases}
g_j + \mathcal{N}(0, \sigma^2) & \text{if } \mathcal{U}(0,1) < p_m \\
g_j & \text{otherwise}
\end{cases}
\end{equation}
where $\sigma$ is the mutation strength, often adapted over generations.

\paragraph{Elitism}
The best $\lfloor e \cdot N \rfloor$ individuals are copied unchanged to the next generation, where $e \in [0, 0.2]$ is the elitism rate.

\paragraph{Convergence Criterion}
The algorithm terminates when:
\begin{itemize}
\item Maximum generations: $t > t_{\max}$
\item Fitness stagnation: Best fitness has not improved for $t_{\text{stale}}$ generations
\item Target fitness: $\text{fit}(\mathbf{x}_{\text{best}}) > 1 - \epsilon_{\text{fit}}$
\end{itemize}

\paragraph{Algorithm}
\begin{enumerate}
\item Initialize population $P(0)$ randomly
\item Evaluate fitness for all $\mathbf{x}_i^{(t)} \in P(t)$
\item Apply elitism: copy top $e \cdot N$ individuals to $P(t+1)$
\item Select parents via tournament selection
\item Apply crossover with probability $p_c$
\item Apply mutation with probability $p_m$
\item Replace population: $P(t) \leftarrow P(t+1)$
\item If convergence not met, $t \leftarrow t+1$, goto 2
\end{enumerate}

\paragraph{Properties}
\begin{itemize}
\item Time complexity: $O(t_{\max} \cdot N \cdot m)$ for $N$ individuals of length $m$
\item Maintains population diversity, reducing local optima entrapment
\item Space complexity: $O(N \cdot |\mathbf{x}|)$ for storing the population
\item Naturally parallelizable across individuals
\end{itemize}

\subsubsection{Constraint Programming (cp)}
A declarative optimization approach that models the problem as a set of variables, domains, and constraints, using propagation and backtracking to find solutions.

\paragraph{Problem Formulation}
A constraint satisfaction problem (CSP) is defined as a tuple:
\begin{equation}
\mathcal{P} = (\mathcal{X}, \mathcal{D}, \mathcal{C})
\end{equation}
where:
\begin{itemize}
\item $\mathcal{X} = \{x_1, x_2, \ldots, x_n\}$ is the set of decision variables
\item $\mathcal{D} = \{D(x_1), D(x_2), \ldots, D(x_n)\}$ is the set of variable domains
\item $\mathcal{C} = \{c_1, c_2, \ldots, c_m\}$ is the set of constraints
\end{itemize}

\paragraph{Domains}
Each variable $x_i$ has an initial domain $D(x_i) \subseteq \mathbb{R}$ (or discrete set). For PCB layout:
\begin{itemize}
\item Position variables: $D(x_i) = [x_{\min}, x_{\max}] \times [y_{\min}, y_{\max}]$
\item Route variables: $D(x_i) = \{\text{path}_1, \text{path}_2, \ldots\}$
\end{itemize}

\paragraph{Constraint Types}
Constraints are classified as:
\begin{itemize}
\item \textbf{Unary}: $c(x_i)$ --- restricts a single variable
\item \textbf{Binary}: $c(x_i, x_j)$ --- relates two variables
\item \textbf{Global}: $c(x_{i_1}, \ldots, x_{i_k})$ --- involves multiple variables (e.g., \texttt{alldiff})
\end{itemize}

\paragraph{Arc Consistency (AC-3)}
For each binary constraint $c(x_i, x_j)$, enforce arc consistency:
\begin{equation}
\forall v_i \in D(x_i), \exists v_j \in D(x_j) \text{ such that } c(v_i, v_j) = \text{true}
\end{equation}
If no such $v_j$ exists, remove $v_i$ from $D(x_i)$. Repeat until no more domains can be reduced.

\paragraph{Propagation}
When a domain is reduced, propagate effects to all connected constraints:
\begin{enumerate}
\item Maintain a worklist $W$ of affected variables
\item For each $x_i \in W$, re-evaluate all constraints involving $x_i$
\item If any domain $D(x_j)$ is reduced, add $x_j$ to $W$
\item Continue until $W$ is empty (fixpoint reached) or a domain becomes empty (failure)
\end{enumerate}

\paragraph{Backtracking Search}
The solver explores the search tree depth-first:
\begin{enumerate}
\item Select unassigned variable $x_i$ (branching heuristic: smallest domain first)
\item For each value $v \in D(x_i)$ (best-first ordering):
\begin{enumerate}
\item Assign $x_i \leftarrow v$
\item Propagate constraints
\item If propagation succeeds (no empty domains), recurse
\item If propagation fails, backtrack
\end{enumerate}
\end{enumerate}

\paragraph{Optimization Extension}
For optimization (minimizing $f(\mathbf{x})$), add a bound constraint:
\begin{equation}
f(\mathbf{x}) < f(\mathbf{x}_{\text{best}})
\end{equation}
This is tightened each time a better solution is found (branch-and-bound).

\paragraph{Convergence Criterion}
\begin{itemize}
\item Solution found: All variables assigned, all constraints satisfied
\item No solution: Search tree exhausted (problem infeasible)
\item Optimality: All nodes pruned, current best is optimal
\end{itemize}

\paragraph{Properties}
\begin{itemize}
\item Time complexity: Exponential in worst case, but propagation prunes large portions of the search tree
\item Guarantees constraint satisfaction (complete when solution exists)
\item Space complexity: $O(n \cdot |D_{\max}| + m)$ for domain storage and constraint graph
\item Deterministic: Same problem instance yields same solution
\end{itemize}

\subsubsection{Gradient Descent (gd)}
A first-order iterative optimization algorithm that finds local minima by moving in the direction of steepest descent of the objective function.

\paragraph{Objective Function}
Assume the objective function $f: \mathbb{R}^n \rightarrow \mathbb{R}$ is differentiable, where $n$ is the number of design variables (component positions, trace parameters, etc.).

\paragraph{Gradient Computation}
The gradient of $f$ at point $\mathbf{x}_k$ is:
\begin{equation}
\nabla f(\mathbf{x}_k) = \left(\frac{\partial f}{\partial x_1}, \frac{\partial f}{\partial x_2}, \ldots, \frac{\partial f}{\partial x_n}\right)^{\top}
\end{equation}
Computed via automatic differentiation or finite differences:
\begin{equation}
\frac{\partial f}{\partial x_i} \approx \frac{f(\mathbf{x}_k + h\mathbf{e}_i) - f(\mathbf{x}_k - h\mathbf{e}_i)}{2h}
\end{equation}
where $h$ is a small perturbation and $\mathbf{e}_i$ is the $i$-th unit vector.

\paragraph{Update Rule}
The basic gradient descent update is:
\begin{equation}
\mathbf{x}_{k+1} = \mathbf{x}_k - \eta_k \nabla f(\mathbf{x}_k)
\end{equation}
where $\eta_k > 0$ is the learning rate at iteration $k$.

\paragraph{Learning Rate Schedules}
\begin{itemize}
\item \textbf{Constant}: $\eta_k = \eta$
\item \textbf{Decay}: $\eta_k = \eta_0 / (1 + \beta k)$
\item \textbf{Exponential}: $\eta_k = \eta_0 \cdot \gamma^k$
\item \textbf{Adaptive}: $\eta_{k,i} = \eta_0 / \sqrt{s_i + \epsilon}$ (per-parameter, AdaGrad-style)
\end{itemize}
where $\eta_0$ is the initial learning rate, $\beta, \gamma$ are decay parameters, $s_i$ accumulates past squared gradients, and $\epsilon$ prevents division by zero.

\paragraph{Momentum Enhancement}
To accelerate convergence and dampen oscillations:
\begin{align}
\mathbf{v}_{k+1} &= \mu \mathbf{v}_k + \eta_k \nabla f(\mathbf{x}_k) \\
\mathbf{x}_{k+1} &= \mathbf{x}_k - \mathbf{v}_{k+1}
\end{align}
where $\mu \in [0, 1)$ is the momentum coefficient (typical: $0.9$).

\paragraph{Constraint Handling}
For constraints $g_j(\mathbf{x}) \leq 0$, use a penalty method:
\begin{equation}
\tilde{f}(\mathbf{x}) = f(\mathbf{x}) + \lambda \sum_{j=1}^{m} \max(0, g_j(\mathbf{x}))^2
\end{equation}
where $\lambda$ is a penalty weight, increased iteratively:
\begin{equation}
\lambda_{k+1} = \min(\lambda_k \cdot \rho, \lambda_{\max}), \quad \rho > 1
\end{equation}

\paragraph{Convergence Criterion}
\begin{itemize}
\item Gradient norm: $\|\nabla f(\mathbf{x}_k)\| < \epsilon_{\text{grad}}$
\item Parameter change: $\|\mathbf{x}_{k+1} - \mathbf{x}_k\| < \epsilon_{\text{step}}$
\item Function change: $|f(\mathbf{x}_{k+1}) - f(\mathbf{x}_k)| < \epsilon_f$
\item Maximum iterations: $k > k_{\max}$
\end{itemize}

\paragraph{Algorithm}
\begin{enumerate}
\item Initialize $\mathbf{x}_0$, $\eta_0$, $\lambda_0$, $k = 0$
\item Compute $\nabla f(\mathbf{x}_k)$ and $\nabla g_j(\mathbf{x}_k)$
\item Update penalty: $\lambda_{k+1} = \min(\lambda_k \cdot \rho, \lambda_{\max})$
\item Compute search direction: $\mathbf{d}_k = -\nabla \tilde{f}(\mathbf{x}_k)$
\item Line search: find $\eta_k$ satisfying Armijo condition
\item Update: $\mathbf{x}_{k+1} = \mathbf{x}_k + \eta_k \mathbf{d}_k$
\item If convergence not met, $k \leftarrow k+1$, goto 2
\end{enumerate}

\paragraph{Properties}
\begin{itemize}
\item Time complexity: $O(k_{\max} \cdot n)$ per iteration for gradient computation
\item Converges to local minima; requires convexity or good initialization for global optimality
\item Space complexity: $O(n)$ for storing the current point and gradient
\item Fast per-iteration, but may require many iterations for convergence
\end{itemize}

\subsubsection{Hybrid (hybrid)}
A multi-strategy approach that combines the strengths of different algorithms in a sequential pipeline, adapting the optimization strategy to the problem characteristics at each phase.

\paragraph{Architecture}
The hybrid solver operates in three sequential phases:
\begin{equation}
\mathbf{x}_{\text{final}} = \text{Phase}_3(\text{Phase}_2(\text{Phase}_1(\mathbf{x}_0)))
\end{equation}

\paragraph{Phase 1: Global Exploration (Simulated Annealing)}
\begin{itemize}
\item Algorithm: Simulated annealing with reduced iterations ($k_{\max}^{(1)} = k_{\max} / 4$)
\item Purpose: Explore the solution space broadly, identify promising regions
\item Output: $\mathbf{x}_1$ with $f(\mathbf{x}_1) < f_{\text{threshold}}$
\end{itemize}
The temperature schedule uses faster cooling ($\alpha = 0.95$) to quickly identify basins of attraction.

\paragraph{Phase 2: Local Refinement (Gradient Descent)}
\begin{itemize}
\item Algorithm: Gradient descent with momentum
\item Purpose: Refine the solution from Phase 1 to a local minimum
\item Input: $\mathbf{x}_1$ from Phase 1
\item Output: $\mathbf{x}_2$ with $\|\nabla f(\mathbf{x}_2)\| < \epsilon_{\text{grad}}$
\end{itemize}
The penalty weight is initialized from Phase 1's final value and continues to increase.

\paragraph{Phase 3: Constraint Satisfaction (Constraint Programming)}
\begin{itemize}
\item Algorithm: Constraint propagation and backtracking
\item Purpose: Verify and enforce all hard constraints
\item Input: $\mathbf{x}_2$ from Phase 2
\item Output: $\mathbf{x}_{\text{final}}$ feasible or infeasible report
\end{itemize}
If $\mathbf{x}_2$ violates hard constraints, CP attempts repair. If repair fails, the solver reports the conflict set.

\paragraph{Adaptive Phase Selection}
Based on problem characteristics, the hybrid solver may skip phases:
\begin{equation}
\text{skip}(\text{Phase}_i) = \begin{cases}
\text{true} & \text{if } n < n_{\text{simple}} \text{ and } i = 1 \\
\text{true} & \text{if } \text{all constraints linear} \text{ and } i = 1 \\
\text{false} & \text{otherwise}
\end{cases}
\end{equation}

\paragraph{Phase Transition Criteria}
Transition from Phase $i$ to Phase $i+1$ occurs when:
\begin{itemize}
\item Improvement rate falls below threshold: $\Delta f / \Delta k < \epsilon_{\text{rate}}$
\item Maximum phase iterations reached: $k^{(i)} > k_{\max}^{(i)}$
\item Solution quality threshold met: $f(\mathbf{x}) < f_{\text{threshold}}^{(i)}$
\end{itemize}

\paragraph{Algorithm}
\begin{enumerate}
\item Initialize $\mathbf{x}_0$, determine phase parameters from problem size $n$
\item \textbf{Phase 1}: Run SA with $k_{\max}^{(1)}$, obtain $\mathbf{x}_1$
\item \textbf{Phase 2}: Run GD from $\mathbf{x}_1$ until convergence, obtain $\mathbf{x}_2$
\item \textbf{Phase 3}: Run CP from $\mathbf{x}_2$, verify/repair constraints
\item Return $\mathbf{x}_{\text{final}}$ or infeasibility report
\end{enumerate}

\paragraph{Properties}
\begin{itemize}
\item Time complexity: Sum of phase complexities, typically $O(k_{\max}^{(1)} \cdot |\mathcal{N}| + k_{\max}^{(2)} \cdot n + \text{CP}_{\text{search}})$
\item Combines global exploration with local exploitation and constraint guarantees
\item Adapts to problem complexity via phase skipping
\item Space complexity: Maximum of individual phase requirements
\end{itemize}

\subsection{Default Algorithm}
If no algorithm is specified, the solver uses simulated annealing (\texttt{sa}) as the default.

\section{Constraint Resolution}
The solver resolves physical constraints according to the following rules:

\subsection{No Conflicts}
If all constraints can be satisfied without conflicts, the solution is complete.

\subsection{Conflicts Between Different Priorities}
When constraints with different priorities conflict:
\begin{itemize}
\item The higher-priority constraint is satisfied
\item The lower-priority constraint is optimized as much as possible without impairing the higher rule
\item A warning is printed to the console
\end{itemize}

\subsection{Conflicts Between Same Priorities}
When constraints with the same priority conflict, behavior depends on the priority level:

\subsubsection{Lower Priorities (-d, -nth, -imp, -vimp)}
\begin{itemize}
\item An error message is printed to the console
\item The program continues execution
\end{itemize}

\subsubsection{Higher Priorities (-oeo, -eo, -ne, -nonneg)}
\begin{itemize}
\item An error message is printed to the console
\item The program does NOT continue execution (halts)
\end{itemize}

\section{Layout Generation Pipeline}
The final layout is generated through a multi-phase pipeline:

\subsection{Phase 1: Constraint Collection}
All physical constraints are collected from the program and organized by type and priority.

\subsection{Phase 2: Board Setup}
The board stackup is defined (layers, thicknesses, materials) and the design area is initialized.

\subsection{Phase 3: Component Placement}
Parts are placed on the board according to:
\begin{enumerate}
\item Alignment constraints (\texttt{phys.align})
\item Proximity constraints (\texttt{phys.prox})
\item Layer constraints (\texttt{phys.layer})
\item Creepage constraints (\texttt{phys.creep})
\end{enumerate}

The solver optimizes placement to satisfy all constraints simultaneously, respecting the priority hierarchy.

\subsection{Phase 4: Trace Routing}
Connections are routed as traces on the board:
\begin{enumerate}
\item Current capacity constraints (\texttt{phys.connectioncurrent}) determine minimum trace width
\item Connection matching constraints (\texttt{phys.connmatch}) ensure length/impedance matching
\item Net shrinkage (\texttt{phys.shrinknet}) minimizes trace length
\end{enumerate}

\subsection{Phase 5: Via Placement}
Vias are placed to connect traces across layers:
\begin{enumerate}
\item Via fill constraints (\texttt{phys.viafill}) populate thermal vias
\item System flags (\texttt{syst.flag.vias}, etc.) control via availability
\end{enumerate}

\subsection{Phase 6: Plane Generation}
Copper planes are generated on specified layers:
\begin{enumerate}
\item Plane constraints (\texttt{phys.mkplane}) create ground and power planes
\item Clearances are maintained around non-connected objects
\end{enumerate}

\subsection{Phase 7: Design Rule Check}
The final layout is verified against all constraints:
\begin{enumerate}
\item All physical constraints are checked
\item Violations are reported with severity based on priority
\item Critical violations (\texttt{-ne}, \texttt{-nonneg}) halt generation
\end{enumerate}

\subsection{Phase 8: File Export}
The verified layout is exported to KiCad format:
\begin{enumerate}
\item \texttt{make.schematic} generates the schematic file (.kicad\_sch)
\item \texttt{make.pcb} generates the PCB file (.kicad\_pcb)
\item \texttt{halt} triggers final export and log file generation
\end{enumerate}

\section{Ordering of Global Set Constant *e}
Inside the \texttt{*e} pseudo-constant (and when converting any set to a list via \texttt{f.list.fromSet()}), element order is determined by lexicographic sorting of each element's \textbf{canonical string representation}.

\subsection{Canonical String Rules}
Each data type produces a deterministic canonical string:
\begin{itemize}
\item \textbf{int} -- decimal string (e.g., \texttt{42} $\rightarrow$ \texttt{"42"})
\item \textbf{real} -- decimal string (e.g., \texttt{3.14} $\rightarrow$ \texttt{"3.14"})
\item \textbf{bool} -- literal (e.g., \texttt{True} $\rightarrow$ \texttt{"True"})
\item \textbf{string} -- quoted (e.g., \texttt{"a"} $\rightarrow$ \texttt{'"a"'})
\item \textbf{list} -- bracket-enclosed, preserving element order (e.g., \texttt{[3, 1, 2]} $\rightarrow$ \texttt{"[3, 1, 2]"})
\item \textbf{set} -- brace-enclosed, elements sorted first (e.g., \texttt{\{Q1, R5, Q2\}} $\rightarrow$ \texttt{"\{Q1, Q2, R5\}"})
\item \textbf{part} -- reference designator (e.g., \texttt{R1} $\rightarrow$ \texttt{"R1"})
\item \textbf{pin} -- parent ref and pad id (e.g., \texttt{U1[3]} $\rightarrow$ \texttt{"U1[3]"})
\item \textbf{connection} -- sorted pin pair in parentheses (e.g., \texttt{(R1[1], U1[+])} $\rightarrow$ \texttt{"(R1[1],U1[+])"})
\item \textbf{net} -- name (e.g., \texttt{Gnd} $\rightarrow$ \texttt{"Gnd"})
\end{itemize}

Elements are then sorted lexicographically by their canonical strings.

\subsection{Example}
\begin{verbatim}
elec.part "U1": opamp(type="LM741"); des=""; 
    fp="Package_DIP:DIP-8_W7.62mm"
elec.net "V-"
elec.connect V- U1[v-]

#c# Canonical strings:
#c# (U1[V-],U1[v-])  -- connection
#c# U1                -- part
#c# U1[+]             -- pin
#c# U1[-]             -- pin
#c# U1[out]           -- pin
#c# U1[v+]            -- pin
#c# U1[v-]            -- pin
#c# V-                -- net
#c#
#c# *e = { (U1[V-],U1[v-]), U1, U1[+], U1[-], 
#c#   U1[out], U1[v+], U1[v-], V- }
\end{verbatim}


